% Figure: shear stress distribution in a circular shaft under torsion. Cross
% section with linearly increasing shear stress from zero at center to tau_max
% at the surface, tau(r) = T r / J.
\begin{figure}[htbp]
\centering
\begin{tikzpicture}[scale=1.0, >=latex]
% the shaft cross-section
\draw[thick] (0,0) circle [radius=2];
\filldraw[enggray] (0,0) circle [radius=1.2pt];
% radius line
\draw[enggray, dashed] (0,0) -- (2,0);
% shear arrows tangent, growing with radius (on the +x radius line, tangent is +y)
\foreach \r in {0.5,1.0,1.5,2.0}{
  \draw[->, engblue, thick] (\r,0) -- (\r,{0.35*\r});
}
% distribution triangle to the right
\begin{scope}[shift={(3.2,0)}]
  \draw[->] (0,-2.2) -- (0,2.2) node[above, font=\scriptsize] {$r$};
  \draw[->] (0,0) -- (2.6,0) node[right, font=\scriptsize] {$\tau$};
  \draw[engblue, very thick] (0,0) -- (2.2,2);
  \draw[engblue, very thick] (0,0) -- (2.2,-2);
  \node[engblue, font=\scriptsize, anchor=west] at (2.2,2) {$\tau_{\max}$};
  \draw[dashed, enggray] (2.2,2) -- (0,2);
  \node[font=\scriptsize, anchor=east] at (0,2) {$R$};
\end{scope}
\node[font=\scriptsize] at (0,-2.6) {circular shaft, torque $T$};
\node[font=\scriptsize] at (3.2,-2.6) {$\tau(r)=Tr/J$};
\end{tikzpicture}
\caption[Shear stress distribution in a shaft under torsion]{Shear
stress in a circular shaft under torque $T$. The shear stress is purely
circumferential and grows linearly with radius, $\tau(r)=Tr/J$, where
$J=\pi d^4/32$ is the polar second moment of area for a solid shaft of
diameter $d$. The stress is zero on the axis and maximum at the surface,
$\tau_{\max}=TR/J=16T/(\pi d^3)$. Because the material near the axis
carries little stress, hollow shafts of the same outer diameter are far
more efficient in torsion per unit weight. Torsion is developed in
Chapter 7 of this volume.}
\label{fig:vol03ch03:torsion-distribution}
\end{figure}
